\section{Integrationstests}
Nach Abschluss der Modultests erfolgt die Zusammenführung der einzelnen Komponenten zum Gesamtsystem.  
Dabei wird insbesondere das Zusammenspiel zwischen Hardware, Sensorik, Netzwerkkommunikation und \gls{led}-Ansteuerung überprüft.  
Bei den Integrationstests werden alle \gls{led}-Stränge gleichzeitig verwendet und auch der Boot-Modus, welcher über die \gls{led}s mehrerer Funktionen läuft, getestet.  
Für die Spannungsversorgung bei den ersten Tests wird ein vorhandenes 25~W Netzteil mit einer Nennspannung von 5~V und einem Nennstrom von 5~A verwendet, welches bereits bei den Modultests eingesetzt wird.  

Die erste Dimensionierung des Netzteils erfolgt auf Grundlage der theoretisch maximal gleichzeitig aktivierten \gls{led}s innerhalb eines Zeitraums von 2~s.  
Die längste mögliche Anzeige aktiviert dabei 104~\gls{led}s.  
Hierbei werden zusätzlich zur Textanzeige auch die \gls{led}s für die maximale Innenraumtemperatur, die maximale Luftfeuchtigkeit, die Mondphase, die Tageszeit sowie die aktuelle Sekunde berücksichtigt.  
Die längste mögliche Hauptanzeige besteht aus 81~\gls{led}s und lautet: „IN REMSECK IST ES FÜNFZEHN IIII NACH SIEBEN MIT SOMMERLICHEN TEMPERATUREN UND BEWÖLKTEM HIMMEL“.  
Zusätzlich sind die 23~\gls{led}s für die aufgezählten Zusatzfunktionen aktiv.  
Für die Berechnung wird, wie in Kapitel~6.2 beschrieben, eine Leistungsaufnahme von 0{,}2~W pro \gls{led} verwendet.  
Daraus ergibt sich eine elektrische Gesamtleistung von

\[
P = 0{,}2~\mathrm{W} \cdot 104 = 20{,}8~\mathrm{W}.
\]

Dies entspricht einem Strom von

\[
I = \frac{20{,}8~\mathrm{W}}{5~\mathrm{V}} = 4{,}16~\mathrm{A}
\]

für die \gls{led}-Beleuchtung.  

Es wird zunächst angenommen, dass die während des Boot-Modus auftretenden Leistungsspitzen durch das Netzteil abgefangen werden können.  
Während der Inbetriebnahme zeigt sich jedoch ein Einbruch der Versorgungsspannung während des Boot-Modus.  
Dadurch startet das System nicht zuverlässig und bleibt zu etwa 75~\% im Boot-Modus hängen.  
Auf Grundlage dieser Beobachtungen erfolgt eine Neuberechnung der erforderlichen Netzteilleistung mit zusätzlicher Leistungsreserve.  

Als Ursache wird identifiziert, dass das Netzteil nicht in der Lage ist, die auftretenden Stromspitzen abzufangen.  
Aus diesem Grund wird der Boot-Modus als Grundlage für die Leistungsberechnung des Systems verwendet.  
Die Berechnung aus Kapitel~6.2 ergibt damit eine maximale Leistungsaufnahme von 51{,}2~W während des Boot-Modus.  

Weitere Tests zeigen, dass ein Netzteil mit einer Leistung von 50~W für den Dauerbetrieb ausreichend dimensioniert ist.  
Das Netzteil stellt sowohl für den Betrieb der maximal 104 gleichzeitig aktiven \gls{led}s als auch für den \gls{esp32c6} und den \gls{dht22} ausreichend Leistungsreserve bereit.  
Zusätzlich zeigen die Tests, dass die während des Boot-Modus auftretenden kurzzeitigen Stromspitzen damit zuverlässig abgefangen werden.  

Nachdem die Spannungsversorgung stabil ist, wird mit den weiteren Integrationstests fortgefahren.  
Zunächst wird die Zeitanzeige durch Vergleich mit der Referenzzeit auf einem Smartphone überprüft, während alle \gls{led}-Stränge gleichzeitig in Betrieb sind.  
Dabei wird insbesondere beobachtet, ob Sekunden- und Minutenwechsel korrekt dargestellt werden und diese in Echtzeit wechseln oder es bei der Verarbeitung zu Verzögerungen kommt.  

Zur Überprüfung aller Minutenanzeigen wird die WordClock über einen Zeitraum von einer Stunde beobachtet.  
Dadurch kann jede mögliche Minutenanzeige mindestens einmal geprüft werden.  
Während dieser Tests wird ein Fehler bei der sprachlichen Zuordnung der Uhrzeit festgestellt.  
Ab dem Bereich „zwanzig nach“ stimmt der Bezug zur folgenden Stunde zunächst nicht korrekt.  
Die Logik zur Berechnung der anzuzeigenden Stunde wird daraufhin angepasst.  

Zusätzlich erfolgt die Überprüfung der Temperatur- und Wetterdarstellung.  
Die angezeigten Wetterzustände werden mit aktuellen Wetterdaten eines Wetterdienstes verglichen.  
Darüber hinaus erfolgt eine subjektive Bewertung der Temperatur- und Wetterklassifikation durch Befragung mehrerer Personen.  
Dabei wird überprüft, ob die gewählten Temperaturbereiche und Wetterbezeichnungen plausibel sind und von Dritten ebenfalls so wahrgenommen werden, ohne dass den befragten Personen die aktuelle Anzeige bekannt ist.  
Es geht dabei vor allem um die Schwellenwerte in der Temperaturanzeige, beispielsweise ab wann von sommerlichen Temperaturen gesprochen wird.  

Die Integrationstests zeigen, dass das Zusammenspiel der einzelnen Komponenten nach der Fehlerkorrektur zuverlässig funktioniert und die Schwellenwerte für Temperatur- und Wetteranzeige gut gewählt sind.  